%! TEX program = pdf_escaped

\documentclass[12pt]{article}

% Packages
\usepackage[a4paper,left=2.5cm,right=2.5cm,top=2.5cm,bottom=2.5cm,footskip=1cm]{geometry}
\usepackage{xcolor}
\usepackage[explicit]{titlesec}
\usepackage[colorlinks=false]{hyperref}
\usepackage[T1]{fontenc}
\usepackage{graphicx}
\usepackage{caption}
\usepackage{listings}
\usepackage{color}
\usepackage{parskip}
\usepackage[ngerman]{babel}
\usepackage{microtype}
\usepackage{svg}
\usepackage{csquotes}
\MakeOuterQuote{"}

% Colors
\definecolor{dkgreen}{rgb}{0,0.6,0}
\definecolor{gray}{rgb}{0.5,0.5,0.5}
\definecolor{mauve}{rgb}{0.58,0,0.82}
\definecolor{lightgray}{rgb}{.9,.9,.9}
\definecolor{darkgray}{rgb}{.4,.4,.4}
\definecolor{purple}{rgb}{0.66, 0.12, 0.82}

% Javascript
\lstdefinelanguage{JavaScript}{
    keywords={typeof, new, let, for, const, continue, of, await, true, false, catch, function, return, null, catch, switch, var, if, in, while, do, else, case, break, try, catch, throw, async },
    keywordstyle=\color{blue}\bfseries,
    ndkeywords={class, export, boolean, default, implements, import, this},
    ndkeywordstyle=\color{darkgray}\bfseries,
    identifierstyle=\color{black},
    sensitive=false,
    comment=[l]{//},
    morecomment=[s]{/*}{*/},
    commentstyle=\color{purple}\ttfamily,
    stringstyle=\color{dkgreen}\ttfamily,
    frame=none,
    morestring=[b]',
    morestring=[b]"
}

% Code
\lstset{frame=tb,
    aboveskip=3mm,
    belowskip=3mm,
    showstringspaces=false,
    columns=flexible,
    basicstyle={\small\ttfamily},
    numbers=left,
    keywordstyle=\color{blue},
    commentstyle=\color{dkgreen},
    stringstyle=\color{mauve},
    breaklines=true,
    breakatwhitespace=true,
    tabsize=4,
    frame=none
}

% Settings
\hypersetup{pdfborder=0 0 0}
\definecolor{FomBlue}{HTML}{00998A}
\linespread{1.20}

% Section heading styling
\titleformat{\subsection} {\bfseries} {\thesubsection} {1em} {#1} [{\titlerule[0.2pt]}]
\titlespacing{\subsection} {0em} {3.3em} {0.8em}

% Title page
\title{%
    \includegraphics[width=5cm]{figures/logo.png}
    \\
    \vspace{1cm}
    \textbf{Sonstige Beteiligung}\\
    \bigskip
    \large Verwendeter Code zur Analyse von Artikeln des Nachrichtensenders \emph{n-tv}
}
\author{Kai Herbst, Manuel Zeh, Henrik Popp}
\date{August 2024}

% Document
\begin{document}
\begin{sloppypar}
	\maketitle
	\thispagestyle{empty}

	\newpage
	\setcounter{page}{1}
	\pagenumbering{Roman}

	\renewcommand{\contentsname}{Inhaltsverzeichnis}
	\tableofcontents

	\newpage
	\setcounter{page}{1}
	\pagenumbering{arabic}

	\section{Einleitung}

	Im Rahmen des Moduls \emph{Analyse semi- \& unstrukturierter Daten} wurden
	die von \emph{n-tv} veröffentlichten Artikel im Zeitraum des 17.04.2024 -
	17.07.2024 gesammelt und nachfolgend analysiert. Das vorliegende Dokument
	enthält den Code zusammen mit dazugehörigen kurzen Beschreibungen, der zur
	Sammlung, Verarbeitung und Analyse der Daten verwendet wurde.

	Der Code ist in die Abschnitte Datenerhebung, Datenvorbereitung,
	Upload-Zeiten-Analyse, Sentiment-Analyse und Topic Modeling unterteilt.

	Jeglicher Code lässt sich ebenfalls in Form von Jupyter Notebooks im
	GitHub-Repository
	\href{https://github.com/okaihe/SS2024-ASUD-Textanalyse/tree/main}{n-tv
		Datenanalyse} finden.


	\newpage
	\section{Datenerhebung}

	Die folgenden Unterkapitel dokumentieren den Code, der zur Sammlung,
	Speicherung und Vorverarbeitung für die späteren Analysen verwendet wurde.

	\subsection{Sammelung und Speicherung der Nachrichtenartikel}

	\hyperref[fig:datenpipeline]{Abbildung \ref{fig:datenpipeline}} zeigt den
	grundlegenden Ablauf, durch den die Nachrichtenartikel gespeichert werden.
	Zunächst wurde in über \emph{AWS Eventbridge} ein Scheduler definiert, der
	täglich um 00:00 Uhr Lambda-Funktion startet, die wiederum auf die von
	\emph{n-tv} angebotenen RSS-Schnittstellen zugreift. In den jeweiligen
	RSS-Feeds für jede Kategorie werden die zuletzt veröffentlichten Artikel
	mit Links und weiteren Metadaten aufgelistet. Über die Links werden die
	Artikel anschließend heruntergeladen und in einem \emph{AWS S3-Bucket}
	gespeichert.

	\begin{figure}[h]
		\centering
		\includesvg[width=12cm]{figures/datenpipeline.svg}
		\caption{Datenpipeline zur Sammeln der \emph{ntv}-Artikel}
		\label{fig:datenpipeline}
	\end{figure}

	Die Lambda-Funktion zum Sammeln der Artikel besteht aus vier Dateien.
	Nachfolgend zu sehen ist die Hauptdatei mit der Einstiegsfunktion:

	\lstinputlisting[language=JavaScript,caption=index.mjs]{../ntv-lambda/index.mjs}

	Die Hauptfunktion (async handler) verwendet wiederum die in den anderen
	drei Dateien \emph{helper.mjs}, \emph{categories.mjs} und
	\emph{s3Helper.mjs} definierten Funktionen auf, die nachfolgend aufgelistet
	werden.

	In der \emph{categories.mjs} werden die abzurufenden Kategorien definiert:

	\lstinputlisting[language=JavaScript,caption=categories.mjs]{../ntv-lambda/categories.mjs}

	Die \emph{s3Helper.mjs}-Datei stellt Funktionen für den Zugriff auf die
	entsprechenden S3-Buckets bereit:

	\lstinputlisting[language=JavaScript,caption=s3Helper.mjs]{../ntv-lambda/s3Helper.mjs}

	In der \emph{helper.mjs}-Datei werden schließlich noch Hilfsfunktionen
	definiert. Unter anderem Funktionen, die einen Großteil des HTML-Codes aus
	den Artikeln bereits im Vorfeld entfernt, um Speicherplatz zu sparen. Dazu
	zählen beispielsweise alle \emph{Script}-Tags, da diese für die weitere
	Bearbeitung nicht benötigt werden. Durch diese Vorgehensweise wird viel
	Speicherplatz und damit auch Kosten eingespart.

	\lstinputlisting[language=JavaScript,caption=helper.mjs]{../ntv-lambda/helper.mjs}

	\subsection{Datenvorbereitung}

	Dieser Code dient dem Erstellen einer CSV-Datei, die anschließend als
	Grundlage für das weitere Vorgehen verwendet wird. Die Funktionen der
	einzelnen Methoden sollten dabei aus den Funktionsnamen hervorgehen.
	Grundsätzlich wird über die Ordner und Dateien des AWS lokal
	heruntergeladenen Buckets iteriert und die Artikel gesammelt. Mit Hilfe von
	\emph{BeautifulSoup} werden Datum und Artikeltext extrahiert und
	gespeichert. Anschließend werden die gesammelten Daten in eine CSV-Datei
	geschrieben. Die CSV-Datei enthält die Spalten date (Datum), time (Uhrzeit),
	category (Kategorie), headline (Überschrift), filename (Dateiname) und text
	(Artikeltext).

	\lstinputlisting[language=Python,caption=Datengenerierung]{./1-1-generation.py}

	\section{Analysevorbereitung}

	Bevor die Daten auf einzelne Aspekte hin analysiert werden, müssen die Daten
	noch weiter aufbereitet und die Sentiment-Analyse und Stichwortgenerierung
	durchgeführt werden.

	\subsection{Vorbereitung der Upload-Analyse}

	Für die Upload-Analyse werden die CSV-Datei um die Spalten upload-hour
	(Stunde des Uploads), weekday (Wochentag) und length (Artikellänge)
	erweitert.

	\lstinputlisting[language=Python,caption=Vorbereitung zur Uploadanalyse]{./1-2-upload.py}

	\subsection{Vorbereitung zur Themenanalyse}

	Für den ersten Ansatz der Themenanalyse müssen Stichwörter zu jedem Artikel
	generiert werden. Dies geschieht mit Hilfe des OpenAI-Sprachmodells
	\emph{GPT-4o mini}, auf das über die OpenAI-REST-API zugegriffen wird und
	das zu jedem Artikel zehn Stichwörter generiert. Die Stichwörter werden
	anschließend der CSV-Datei hinzugefügt:

	\lstinputlisting[language=Python,caption=Vorbereitung zur Stichwortanalyse,inputencoding=utf8]{./1-3-keywords.py}

	\subsection{Durchführung Sentiment-Analyse}

	Um die Sentiments der Artikel genauer untersuchen zu können, müssen diese
	zunächst identifiziert werden. Dazu wird die CSV-Datei folgende Spalten
    erweitert:
    \begin{itemize}
        \item \texttt{sentiment\_headline}
	    \item \texttt{sentiment\_text}
        \item \texttt{sentiment\_prop\_headline\_positive}
	    \item \texttt{sentiment\_prop\_headline\_neutral}
        \item \texttt{sentiment\_prop\_headline\_negative}
	    \item \texttt{sentiment\_prop\_text\_positive}
        \item \texttt{sentiment\_prop\_text\_neutral}
	    \item \texttt{sentiment\_prop\_text\_negative}
    \end{itemize}

	\lstinputlisting[language=Python,caption=Vorbereitung zur Sentiment-Untersuchung,inputencoding=utf8]{./1-4-sentiment.py}

	\section{Analyse der Daten}

	\subsection{Analyse des Uploadverhaltens}

	Im nachfolgenden sind die Funktionen, die zur Erstellung der Diagramme
	verwendet wurden. Diese sollten in einzelnen Zellen in einem
	Jupyter-Notebook ausgeführt werden, um die Ergebnisse direkt sehen zu
	können.

	\lstinputlisting[language=Python,caption=Analyse der Uploadzeiten,inputencoding=utf8]{./3-1-upload-analyse.py}

	\subsection{Analyse der Sentiments}

	Der im Anschluss aufgeführte Code enthält die Funktionen zur Generierung von
	Visualisierungen zu den Sentiments. Auch diese sollten vorzugsweise in einem
	Jupyter Notebook in einzelnen Zellen ausgeführt werden, um die Diagramme
	sehen zu können.

	\lstinputlisting[language=Python,caption=Analyse der Sentiments,inputencoding=utf8]{./3-2-sentiment-analysis.py}

	\subsection{Topic-Analyse über Stichwörter}

	Ziel der Analyse über die Stichwörter ist die Generierung einer Karte, die
	die Topics als Netzwerk zeigt. Dazu wird eine Graphdatei aus den Daten
	erstellt, die anschließend in Gephi importiert werden kann. Die Gephi-Datei
	ist ebenso dem Upload angefügt. Die Stichwörter werden dazu zunächst
	auseinandergenommen und in einem neuen Dataframe gespeichert. Anschließend
	wird über die Zeilen iteriert und Informationen über Häufigkeite der
	einzelnen Stichwörter und Häufigkeit über Kombinationen von Stichwörtern
	gesammelt. Diese Informationen werden anschließend \emph{networkx}
	übergeben, der daraus ein GraphML-Datei erzeugt, die in Gephi importiert
	werden kann.

	\lstinputlisting[language=Python,caption=Erstellung der Graph-Datei,inputencoding=utf8]{./3-3-keywords.py}

	\subsection{Topic-Analyse über Topic Modeling}

	Als zweiter Ansatz zur Untersuchung der Themen wurden \emph{BERTopic}
	verwendet. Bertopic bietet neben der Extrahierung der Themen auch zahlreiche
	Funktionen zur Visualisierung. Wie in den anderen Abschnitten sollte auch
	der Code zur Visualisierung in einem Jupyter Notebook in einzelnen Zellen
	aufgerufen werden. So können die Ergebnisse direkt betrachtet und interaktiv
	verwendet werden.

	\lstinputlisting[language=Python,caption=Erstellung der Graph-Datei,inputencoding=utf8]{./3-4-topic-modeling.py}

    \newpage
	\section{Fazit}

	Dank moderner Algorithmen und zahlreicher Python-Packages lassen sich solche
	Analyseprojekte heutzutage auch mit relativ wenigen Vorkenntnissen im
	Themenbereich gut umsetzen. Gerade das Package \emph{Matplotlib} stellt bei
	der Erstellung Diagrammen bzw. der Visualisierung von Ergebnissen eine große
	Hilfe dar. Dank der modernen Entwicklung im Bereich KI lassen sich auch
	nachträglich zu fremden Texten Stichwörter generieren, was bei der
	Erforschung der Themen und der Erstellung des Netzwerks stark geholfen hat.
	\emph{BERTopic} hat sich ebenfalls als gute Wahl zur Erforschung der Themen
	erwiesen, da es gleichzeitig zahlreiche Visualisierungsmöglichkeiten bietet.

\end{sloppypar}
\end{document}
